\section{Proof of Theorem \ref{thm:loss_transformation}} \label{sec:thm_proof}

\begin{proof}
    
We know that the elements of a vector $u$ that is centralized and normalized (using the $\mathcal{N}$ operator mentioned above) have a variance of $1/n$, since:
\begin{equation}
    \begin{aligned}
        \var\brackets{u}
        & = \expectation{i}{u_i^2} - \expectation{i}{u_i}^2 \\
        & = \frac{1}{n}\|u\|^2 - 0 \\
        & = \frac{1}{n}.
    \end{aligned}
\end{equation}

Using the original definition of Pearson's correlation coefficient, we have:
\begin{equation}
    \begin{aligned}
        \sigma(u, v)
        & = \frac{\cov\brackets{u, v}}{\sqrt{\var\brackets{u} \cdot \var\brackets{v}}} \\
        & = \expectation{i}{
            \frac{(u_i - \bar{u})}{\sqrt{\var\brackets{u}}}
            \cdot
            \frac{(v_i - \bar{v})}{\sqrt{\var\brackets{v}}}
        } \\
        & = \expectation{i}{
            \frac{\brackets{\mathcal{N}(u)}_i}{\sqrt{1/n}}
            \cdot
            \frac{\brackets{\mathcal{N}(v)}_i}{\sqrt{1/n}}
        } \\
        & = n\expectation{i}{\brackets{\mathcal{N}(u)}_i \brackets{\mathcal{N}(v)}_i} \\
        & = \angles{\mathcal{N}(u), \mathcal{N}(v)}.
    \end{aligned}
\end{equation}
That is to say, the Pearson's correlation coefficient between two vectors equals the inner product of the two vectors centralized and normalized.

Therefore the theorem can be proved as follows. Recall that $f_i(x_t)$ and $y_t$ are normalized.
\begin{equation}
    \begin{aligned}
        n\mathcal{L}(w)
        & = \frac{1}{T} \sum_{t=1}^T \|z_t - y_t\|_2^2 \\
        & = \expectation{t}{\|z_t\|^2 - 2\angles{z_t, y_t} + \|y_t\|^2} \\
        & = \expectation{t}{
            \Verts{\sum_{i=1}^k w_if_i(X_t)}^2
            - 2 \angles{\sum_{i=1}^k w_if_i(X_t), y_t}
            + 1
        } \\
        & = \mathbb{E}_t\left[ \sum_{i=1}^k\sum_{j=1}^k w_iw_j \sigma(f_i(X_t), f_j(X_t))\right. \\
        & \qquad \left. -\ 2\sum_{i=1}^k w_i\angles{f_i(X_t), y_t} + 1 \right] \\
        & = \sum_{i=1}^k\sum_{j=1}^k w_iw_j \bar{\sigma}(f_i(X), f_j(X)) - 2\sum_{i=1}^k w_i \bar{\sigma}_y(f_i) + 1.
    \end{aligned}
\end{equation}

\end{proof}
